\documentclass[letterpaper]{article}
\usepackage[margin=1in]{geometry}
\usepackage{amsmath,amssymb}
\usepackage{booktabs}
\usepackage{multirow}
\usepackage{graphicx}
\usepackage{xcolor}
\usepackage{hyperref}
\usepackage[numbers]{natbib}
\usepackage{tabularx}

\newcommand{\ours}{Meta-Brain}
\newcommand{\pp}{\,\text{pp}}

\title{VLA Internals Already Know What Went Wrong:\\Diagnostic Probing and the Limits of\\Perturbation-Routed LoRA Correction}

\author{
  Anonymous Authors\\
  CoRL 2026 Submission
}

\date{}

\begin{document}
\maketitle

%=====================================================================
\begin{abstract}
Vision-Language-Action (VLA) models achieve high success rates on standard manipulation benchmarks but degrade sharply under environmental perturbations. We investigate whether VLA internal representations encode sufficient information about perturbation type to enable targeted correction via expert routing. We report two findings and one negative result. \textbf{Finding 1}: A linear probe on a single layer of a frozen pi0 backbone (PaliGemma layer~6, first token) classifies perturbation type---visual, spatial, or physical---with near-perfect accuracy (3-class F1\,$=$\,0.993; K\,$=$\,1 F1\,$=$\,0.898), and this phenomenon replicates on OpenVLA (F1\,$=$\,0.952). Layer placement analysis reveals a consistent depth ordering: visual perturbation signal peaks in early layers, spatial in middle layers, and physical in late layers. \textbf{Finding 2}: LoRA adaptation of any form substantially improves robustness over frozen baselines (+10.4$\pp$ for routed experts, +9.7$\pp$ for a single rank-48 LoRA). \textbf{Negative result}: Perturbation-type-specific routing does not outperform a parameter-matched universal LoRA. On LIBERO-Plus, our routed system (\ours{}, 64.0\%) is statistically indistinguishable from a rank-48 universal LoRA (63.3\%, $p$\,$=$\,0.42) and within 1.6$\pp$ of random routing (62.4\%, $p$\,$=$\,0.18). An oracle routing upper bound (69.9\%) reveals a 5.9$\pp$ gap, suggesting that routing could become valuable if experts learn more specialized corrections---an open problem we characterize but do not solve. We additionally find that routing introduces a measurable $-4.2\pp$ degradation on unperturbed tasks. We release our diagnostic toolkit and benchmark results as a resource for future VLA robustness research.
\end{abstract}

%=====================================================================
\section{Introduction}
\label{sec:intro}

VLA models such as pi0 and OpenVLA achieve success rates above 90\% on standard manipulation benchmarks like LIBERO, but their performance degrades substantially under environmental perturbations. On LIBERO-Plus---a robustness benchmark introducing visual, spatial, and physical perturbations---a frozen pi0 achieves only 53.6\%, a drop of nearly 40 percentage points.

Existing robustness methods either apply perturbation-agnostic corrections (RobustVLA, VLA-GSE) or target a single perturbation modality (FTM for visual perturbations). None attempt to \emph{diagnose} the perturbation type and route to a specialized corrector. This paper asks: \emph{do VLA internals encode enough perturbation-type information to enable such targeted correction, and does targeted correction actually help?}

We answer the first question affirmatively. A linear probe on PaliGemma layer~6 of a frozen pi0 achieves 3-class F1 of 0.993 for distinguishing visual, spatial, and physical perturbations---and this finding replicates on OpenVLA (F1\,$=$\,0.952). The perturbation signal follows a consistent depth ordering across architectures: visual information peaks early, spatial in mid-layers, and physical late. These findings establish that VLA representations are a rich, untapped diagnostic resource.

For the second question, we built \ours{}: a system that routes a frozen VLA's activations through a Perturbation-Aware Router (PAR) to a bank of perturbation-type-specific LoRA experts. Despite the near-perfect routing accuracy, the routed system (64.0\% on LIBERO-Plus) does not statistically outperform a single rank-48 universal LoRA (63.3\%, $p=0.42$) that uses the same parameter budget without routing. Random routing achieves 62.4\%, confirming that the multi-expert benefit is primarily an ensemble/capacity effect rather than type-specific specialization.

Rather than burying this negative result, we present it as the paper's central empirical lesson: \emph{diagnosing perturbation type is easy; exploiting that diagnosis for correction is hard}. An oracle routing experiment (69.9\%) reveals a 5.9$\pp$ gap over learned routing, indicating unrealized potential that current expert training methods cannot unlock. We frame this gap as an open problem and discuss concrete directions for future work.

Our contributions are:
\begin{enumerate}
    \item A mechanistic finding that VLA internal representations encode perturbation type at single-layer, single-step granularity, with a consistent depth ordering (V$\to$early, S$\to$mid, P$\to$late), replicated across two VLA architectures.
    \item A comprehensive ablation study establishing that perturbation-type routing does not outperform parameter-matched universal adaptation under current expert training, presented as an honest negative result with statistical rigor.
    \item Characterization of the oracle gap as a concrete open problem, with rank-sweep evidence suggesting that expert specialization may emerge at higher capacity.
\end{enumerate}

%=====================================================================
\section{Related Work}
\label{sec:related}

\paragraph{VLA robustness.}
Several methods address VLA fragility under distribution shift. FTM applies fine-tuning with masking for visual perturbations, achieving strong results (71.2\% on visual categories) but cannot handle spatial or physical perturbations. RobustVLA uses uniform adversarial training. VLA-GSE routes tokens by capacity needs using a Mixture-of-Experts architecture (114M parameters), achieving 81.2\% on LIBERO-Plus with a Qwen3-VL-4B backbone. None of these methods diagnose perturbation type to inform the correction strategy.

\paragraph{Mixture-of-Experts and LoRA routing.}
MoIRA and CORAL route among LoRA experts for multi-task adaptation, but their routing criteria are task identity rather than perturbation type. LoRA-SP adapts rank allocation to task complexity. Our work differs in routing by diagnosed environmental perturbation type---a dimension orthogonal to task or capacity routing.

\paragraph{Representation probing in foundation models.}
Linear probing of internal representations is well-established for understanding what models encode. SAFE demonstrates that VLA activations can detect perturbation presence. We extend this to perturbation \emph{type} classification and show the signal is accessible at a single layer and single timestep, with a depth-ordered structure across perturbation categories.

%=====================================================================
\section{Method}
\label{sec:method}

\subsection{Phase 0: Perturbation Diagnostic Probing}
\label{sec:phase0}

Given a frozen VLA backbone $f_\theta$ with $L$ layers, we extract intermediate activations $\mathbf{h}^{(l)} \in \mathbb{R}^d$ at each layer $l$ for episodes executed under known perturbation types $y \in \{\text{Visual}, \text{Spatial}, \text{Physical}\}$. We train a linear probe $g_l: \mathbb{R}^d \to \mathbb{R}^3$ at each layer to predict $y$ from $\mathbf{h}^{(l)}$.

For pi0 (PaliGemma backbone, 18 layers), the optimal probing layer is layer~6 with first-token pooling, achieving 3-class F1\,$=$\,0.993 on held-out episodes. Critically, single-step (K\,$=$\,1) classification achieves F1\,$=$\,0.898 at layer~6 and F1\,$=$\,1.000 with first-token pooling, establishing that perturbation type is decodable from a single forward pass without temporal aggregation.

\subsection{Phase 1: \ours{} Architecture}
\label{sec:architecture}

\ours{} consists of three components applied on top of a frozen VLA:

\paragraph{Perturbation-Aware Router (PAR).}
A 2-layer MLP ($\sim$0.53M parameters) operating on layer~6 PaliGemma activations, producing a soft routing distribution $\mathbf{w} = \text{softmax}(\text{MLP}(\mathbf{h}^{(6)})) \in \Delta^2$ over three expert indices. Online classification accuracy is 97.8\%.

\paragraph{LoRA Expert Bank.}
Three rank-16 LoRA adapters ($\sim$3.7M parameters each) applied to the attention layers of the frozen backbone: LoRA-V (trained on visual perturbation data), LoRA-S (spatial), and LoRA-P (physical). Each expert is first trained independently on its perturbation type, then frozen during router fine-tuning.

\paragraph{Soft Merging.}
The corrected output is computed as the weighted combination: $\Delta\theta = \sum_i w_i \Delta\theta_i$, where $\Delta\theta_i$ are the LoRA expert parameters and $w_i$ the routing weights.

Total trainable parameters: $\sim$12.1M (three experts $\times$ 3.7M + router 0.53M).

%=====================================================================
\section{Experiments}
\label{sec:experiments}

We evaluate on LIBERO-Plus (7 perturbation subtypes across 3 categories) and standard LIBERO (unperturbed). All results use 5 random seeds unless otherwise noted.

%---------------------------------------------------------------------
\subsection{Expert Cross-Evaluation: Limited Type Specificity}
\label{sec:cross_eval}

We first evaluate each perturbation-specific LoRA expert on all perturbation types to measure type specificity (Table~\ref{tab:cross_eval}).

\begin{table}[h]
\centering
\caption{Expert $\times$ Perturbation Type cross-evaluation (success rate \%). Cross-type performance drops average only 4.7$\pp$, indicating experts learn largely interchangeable corrections.}
\label{tab:cross_eval}
\small
\begin{tabular}{lcccc}
\toprule
\textbf{Method} & \textbf{Visual} & \textbf{Spatial} & \textbf{Physical} & \textbf{Overall} \\
\midrule
LoRA-V (rank 16) & 67.0 $\pm$ 2.1 & 58.0 $\pm$ 2.8 & 49.5 $\pm$ 3.2 & 60.8 $\pm$ 1.9 \\
LoRA-S (rank 16) & 62.5 $\pm$ 2.4 & 62.0 $\pm$ 2.6 & 50.0 $\pm$ 3.0 & 59.7 $\pm$ 1.8 \\
LoRA-P (rank 16) & 61.0 $\pm$ 2.5 & 57.5 $\pm$ 2.9 & 55.5 $\pm$ 3.1 & 59.2 $\pm$ 2.0 \\
Universal (rank 16) & 61.5 $\pm$ 2.3 & 56.0 $\pm$ 2.7 & 49.0 $\pm$ 3.3 & 58.5 $\pm$ 1.7 \\
Frozen pi0 & 51.2 $\pm$ 1.8 & 57.4 $\pm$ 2.5 & 54.8 $\pm$ 2.9 & 53.6 $\pm$ 1.5 \\
\bottomrule
\end{tabular}
\end{table}

The cross-type drop when applying a mismatched expert averages only 4.7$\pp$ (Table~\ref{tab:cross_drop}). LoRA-S applied to visual perturbations still reaches 62.5\%, only 4.5$\pp$ below the matched LoRA-V. This small drop indicates that rank-16 LoRA experts learn largely interchangeable corrections rather than type-specific repair functions.

\begin{table}[h]
\centering
\caption{Cross-type mismatch drop. Gaps of 4--6$\pp$ are far below the $>$10$\pp$ threshold that would justify routing overhead.}
\label{tab:cross_drop}
\small
\begin{tabular}{lccc}
\toprule
\textbf{Perturbation Type} & \textbf{Matched Expert} & \textbf{Best Mismatched} & \textbf{Drop ($\pp$)} \\
\midrule
Visual & LoRA-V: 67.0 & LoRA-S: 62.5 & 4.5 \\
Spatial & LoRA-S: 62.0 & LoRA-V: 58.0 & 4.0 \\
Physical & LoRA-P: 55.5 & LoRA-S: 50.0 & 5.5 \\
\midrule
\multicolumn{3}{l}{\textbf{Average drop}} & \textbf{4.7} \\
\bottomrule
\end{tabular}
\end{table}

%---------------------------------------------------------------------
\subsection{Main Result: Ablation Ladder}
\label{sec:ablation_ladder}

Table~\ref{tab:ablation} presents the central ablation ladder comparing all methods on LIBERO-Plus.

\begin{table}[h]
\centering
\caption{Complete ablation ladder on LIBERO-Plus (5 seeds). \ours{} does not statistically outperform a parameter-matched rank-48 universal LoRA ($p=0.42$) or random routing ($p=0.18$). The oracle gap (5.9$\pp$) indicates unrealized routing potential.}
\label{tab:ablation}
\small
\begin{tabular}{llccccc}
\toprule
\textbf{Method} & \textbf{Params} & \textbf{Overall (\%)} & \textbf{Visual} & \textbf{Spatial} & \textbf{Physical} & $\boldsymbol{\Delta}$ \textbf{vs Frozen} \\
\midrule
Frozen pi0 & 0 & 53.6 $\pm$ 1.5 & 51.2 & 57.4 & 54.8 & --- \\
Universal r16 & 3.7M & 58.5 $\pm$ 1.7 & 61.5 & 56.0 & 49.0 & +4.9 \\
Random routing & 11.6M & 62.4 $\pm$ 2.2 & 65.5 & 60.0 & 52.5 & +8.8 \\
Universal r48 & 12.0M & 63.3 $\pm$ 1.9 & 66.5 & 61.0 & 53.0 & +9.7 \\
\ours{} (PAR) & 12.1M & 64.0 $\pm$ 1.6 & 67.0 & 61.5 & 54.0 & +10.4 \\
\ours{} + CM & 12.1M+ & 64.7 $\pm$ 1.7 & 67.6 & 62.0 & 54.6 & +11.1 \\
Oracle routing & 11.6M & 69.9 $\pm$ 1.4 & 73.0 & 68.0 & 59.5 & +16.3 \\
\bottomrule
\end{tabular}
\end{table}

Three observations are central:

\begin{enumerate}
    \item \textbf{Routing $\approx$ scaling.} \ours{} (64.0\%) vs.\ universal r48 (63.3\%): $\Delta = 0.7\pp$, $p = 0.42$. A single LoRA with matched parameter count achieves equivalent performance without any routing mechanism.
    \item \textbf{Routing $\approx$ random.} \ours{} (64.0\%) vs.\ random routing (62.4\%): $\Delta = 1.6\pp$, $p = 0.18$. Random expert assignment captures most of the multi-expert benefit, indicating ensemble/capacity effects dominate over type-specific routing.
    \item \textbf{Oracle gap.} Oracle routing (69.9\%) vs.\ \ours{} (64.0\%): $\Delta = 5.9\pp$, $p = 0.004$. Perfect routing significantly outperforms learned routing, revealing unrealized potential that current expert training cannot exploit.
\end{enumerate}

Table~\ref{tab:significance} reports full statistical tests.

\begin{table}[h]
\centering
\caption{Statistical significance tests (paired $t$-test, 5 seeds).}
\label{tab:significance}
\small
\begin{tabular}{lcccc}
\toprule
\textbf{Comparison} & $\boldsymbol{\Delta}$ \textbf{($\pp$)} & \textbf{$t$-statistic} & \textbf{$p$-value} & \textbf{Sig.?} \\
\midrule
\ours{} vs Frozen & +10.4 & 8.12 & 0.001 & Yes \\
\ours{} vs Universal r16 & +5.5 & 4.31 & 0.013 & Yes \\
\ours{} vs Universal r48 & +0.7 & 0.88 & 0.42 & No \\
\ours{} vs Random & +1.6 & 1.56 & 0.18 & No \\
Oracle vs \ours{} & +5.9 & 5.87 & 0.004 & Yes \\
Random vs Universal r16 & +3.9 & 3.02 & 0.039 & Yes \\
\bottomrule
\end{tabular}
\end{table}

%---------------------------------------------------------------------
\subsection{Parameter-Matched Control: Universal Rank-48}
\label{sec:r48}

The existential control for type-specific routing is a single universal LoRA with rank~48 ($\sim$12.0M parameters), matching \ours{}'s total parameter budget. Table~\ref{tab:r48} shows the comparison.

\begin{table}[h]
\centering
\caption{Parameter-matched comparison. Universal r48 achieves equivalent performance to \ours{} with a simpler architecture and no routing overhead.}
\label{tab:r48}
\small
\begin{tabular}{lccccl}
\toprule
\textbf{Metric} & \textbf{Universal r48} & \textbf{\ours{}} & $\boldsymbol{\Delta}$ \\
\midrule
Params (trainable) & 12.0M & 12.1M & --- \\
Overall (\%) & 63.3 $\pm$ 1.9 & 64.0 $\pm$ 1.6 & +0.7 \\
Visual (\%) & 66.5 & 67.0 & +0.5 \\
Spatial (\%) & 61.0 & 61.5 & +0.5 \\
Physical (\%) & 53.0 & 54.0 & +1.0 \\
Architecture & Single LoRA & Router + 3 LoRA & --- \\
Inference overhead & None & Router fwd pass & --- \\
\bottomrule
\end{tabular}
\end{table}

The 0.7$\pp$ gap is not statistically significant ($p = 0.42$). This demonstrates that at the 12M parameter budget, allocating capacity to a single universal adapter is as effective as distributing it across three specialized experts with learned routing. The additional architectural complexity of the router and soft-merge mechanism provides no measurable benefit.

%---------------------------------------------------------------------
\subsection{\ours{} Per-Perturbation Breakdown and Routing Analysis}
\label{sec:meta_brain_detail}

Table~\ref{tab:perturbation} shows \ours{}'s per-perturbation results on LIBERO-Plus.

\begin{table}[h]
\centering
\caption{\ours{} LIBERO-Plus per-perturbation breakdown.}
\label{tab:perturbation}
\small
\begin{tabular}{llccc}
\toprule
\textbf{Perturbation} & \textbf{Type} & \textbf{\ours{} (\%)} & \textbf{Frozen pi0 (\%)} & $\boldsymbol{\Delta}$ \textbf{($\pp$)} \\
\midrule
Camera & Visual & 69.5 $\pm$ 2.4 & 49.0 $\pm$ 2.0 & +20.5 \\
Light & Visual & 71.0 $\pm$ 2.2 & 55.2 $\pm$ 1.9 & +15.8 \\
Background & Visual & 64.0 $\pm$ 2.6 & 57.0 $\pm$ 2.2 & +7.0 \\
Noise & Visual & 63.5 $\pm$ 2.7 & 47.6 $\pm$ 2.1 & +15.9 \\
Layout & Spatial & 61.5 $\pm$ 2.8 & 57.4 $\pm$ 2.5 & +4.1 \\
Robot & Physical & 54.0 $\pm$ 3.2 & 54.8 $\pm$ 2.9 & $-$0.8 \\
\midrule
Overall & --- & 64.0 $\pm$ 1.6 & 53.6 $\pm$ 1.5 & +10.4 \\
\bottomrule
\end{tabular}
\end{table}

\ours{} improves most on visual perturbations (camera: +20.5$\pp$, noise: +15.9$\pp$), modestly on spatial (+4.1$\pp$), and \emph{slightly degrades} on physical ($-0.8\pp$). Physical perturbations remain the hardest category: LoRA-P's peak performance is only 55.5\%, suggesting that rank-16 LoRA cannot learn effective corrections for changes in contact dynamics and force profiles.

\paragraph{Routing weight analysis.}
Table~\ref{tab:routing_weights} reveals that the PAR router's soft routing weights do not converge to near-one-hot distributions, despite 97.8\% classification accuracy. The system effectively learns a soft ensemble rather than type-specific routing.

\begin{table}[h]
\centering
\caption{PAR soft routing weight matrix. Entropy values near the maximum ($\log 3 = 1.58$) indicate near-uniform distributions rather than type-specific selection.}
\label{tab:routing_weights}
\small
\begin{tabular}{lcccc}
\toprule
\textbf{True Type} & $\boldsymbol{\to}$ \textbf{LoRA-V} & $\boldsymbol{\to}$ \textbf{LoRA-S} & $\boldsymbol{\to}$ \textbf{LoRA-P} & \textbf{Entropy} \\
\midrule
Visual & 0.58 & 0.24 & 0.18 & 1.32 \\
Spatial & 0.28 & 0.52 & 0.20 & 1.41 \\
Physical & 0.22 & 0.30 & 0.48 & 1.48 \\
\bottomrule
\end{tabular}
\end{table}

Even for visual perturbations---the category where routing should be most clear-cut---only 58\% of weight goes to LoRA-V, with 42\% dispersed to the other experts. Physical routing is worst: only 48\% to LoRA-P, with 30\% to LoRA-S. The end-to-end optimization discovers that mixing experts outperforms strict selection, confirming that expert corrections are largely interchangeable.

%---------------------------------------------------------------------
\subsection{Oracle Routing: Quantifying the Ceiling}
\label{sec:oracle}

Replacing the PAR router with ground-truth perturbation labels (hard one-hot routing) yields 69.9\% overall (Table~\ref{tab:oracle}).

\begin{table}[h]
\centering
\caption{Oracle vs PAR routing comparison. The 5.9$\pp$ gap is significant ($p=0.004$).}
\label{tab:oracle}
\small
\begin{tabular}{lcccc}
\toprule
\textbf{Method} & \textbf{Visual} & \textbf{Spatial} & \textbf{Physical} & \textbf{Overall} \\
\midrule
Oracle routing & 73.0 $\pm$ 1.8 & 68.0 $\pm$ 2.3 & 59.5 $\pm$ 2.7 & 69.9 $\pm$ 1.4 \\
\ours{} (PAR) & 67.0 $\pm$ 2.1 & 61.5 $\pm$ 2.8 & 54.0 $\pm$ 3.2 & 64.0 $\pm$ 1.6 \\
Gap ($\pp$) & +6.0 & +6.5 & +5.5 & +5.9 \\
\bottomrule
\end{tabular}
\end{table}

The oracle gap is \emph{not} due to imperfect routing accuracy (PAR achieves 97.8\% accuracy). Rather, the gap arises from the structural difference between hard one-hot and soft routing: when mismatched expert weights are exactly zero, the correction signal is purer and avoids inter-expert interference. Oracle routing also significantly outperforms universal r48 (69.9\% vs 63.3\%, +6.6$\pp$), confirming that type-specific routing has value \emph{in principle}---but only under the idealized condition of hard selection. This gap constitutes a concrete open problem: training experts that are specialized enough to benefit from hard routing.

%---------------------------------------------------------------------
\subsection{Comparison with FTM}
\label{sec:ftm}

Table~\ref{tab:ftm} compares \ours{} against a reproduced FTM baseline.

\begin{table}[h]
\centering
\caption{FTM vs \ours{}. FTM dominates on visual perturbations with $250\times$ fewer parameters, but cannot address spatial or physical perturbations.}
\label{tab:ftm}
\small
\begin{tabular}{lccccc}
\toprule
\textbf{Method} & \textbf{Params} & \textbf{Visual} & \textbf{Spatial} & \textbf{Physical} & \textbf{Overall} \\
\midrule
FTM (reproduced) & 4K & 71.2 & 55.0 & 52.0 & 62.8 \\
\ours{} & 12.1M & 67.0 & 61.5 & 54.0 & 64.0 \\
LoRA-V only & 3.7M & 67.0 & 58.0 & 49.5 & 60.8 \\
\bottomrule
\end{tabular}
\end{table}

FTM achieves 71.2\% on visual perturbations with only 4K parameters, exceeding \ours{} by 4.2$\pp$ on this category. \ours{}'s advantage is breadth (+6.5$\pp$ on spatial, +2.0$\pp$ on physical vs.\ FTM), but this breadth advantage is also achieved by the simpler rank-48 universal LoRA. This comparison reinforces our finding: specialized methods (FTM for visual) outperform routing-based approaches on their target domain, while general robustness is better served by capacity scaling than type-aware routing.

%---------------------------------------------------------------------
\subsection{Parameter Efficiency Analysis}
\label{sec:param_efficiency}

Table~\ref{tab:param_eff} reports gain-over-frozen-baseline normalized by trainable parameters.

\begin{table}[h]
\centering
\caption{Parameter efficiency comparison. Routing does not improve efficiency over universal adaptation at matched parameter budgets.}
\label{tab:param_eff}
\small
\begin{tabular}{lcccc}
\toprule
\textbf{Method} & \textbf{Params} & \textbf{Gain ($\pp$)} & \textbf{Gain/M params} & \textbf{Coverage} \\
\midrule
FTM & 4K & +20.0 (visual) & 4883.0 & Visual only \\
Universal r16 & 3.7M & +4.9 & 1.32 & All types \\
\ours{} & 12.1M & +10.4 & 0.86 & All types \\
Universal r48 & 12.0M & +9.7 & 0.81 & All types \\
VLA-GSE$^*$ & 114M & +12.0 & 0.105 & All types \\
\bottomrule
\multicolumn{5}{l}{\footnotesize $^*$Different backbone (Qwen3-VL-4B); indirect comparison only.}
\end{tabular}
\end{table}

\ours{} (0.86 gain/M) and universal r48 (0.81 gain/M) are nearly identical in parameter efficiency. Routing does not yield a parameter efficiency advantage, as the router's parameters ($\sim$0.5M) contribute no marginal performance gain. FTM's extreme efficiency on visual perturbations (4883 gain/M) highlights that domain-specific methods dominate when the perturbation type is known.

%---------------------------------------------------------------------
\subsection{Correction Memory Ablation}
\label{sec:cm}

An optional Correction Memory (CM) module that modulates routing weights via exemplar retrieval adds 0.7$\pp$ (64.0\%$\to$64.7\%), which is not significant ($p = 0.35$). The learned routing modulation weight ($\alpha = 0.12$) confirms that end-to-end training effectively suppresses CM's influence. We do not include CM in the main system.

%---------------------------------------------------------------------
\subsection{Router Architecture Comparison}
\label{sec:router_arch}

Table~\ref{tab:router_arch} compares three PAR router architectures.

\begin{table}[h]
\centering
\caption{Router architecture comparison. All three achieve $>$98\% classification accuracy, but downstream differences are negligible (0.6$\pp$ range).}
\label{tab:router_arch}
\small
\begin{tabular}{lcccc}
\toprule
\textbf{Router} & \textbf{Params} & \textbf{Routing Acc (\%)} & \textbf{Downstream (\%)} & \textbf{Latency (ms)} \\
\midrule
2-layer MLP & 0.53M & 99.6 & 64.0 $\pm$ 1.6 & 0.3 \\
Linear & 0.02M & 98.1 & 63.4 $\pm$ 1.8 & 0.1 \\
Attention-based & 1.2M & 99.3 & 63.8 $\pm$ 1.7 & 0.8 \\
\bottomrule
\end{tabular}
\end{table}

Even a linear classifier (98.1\% accuracy, 0.02M parameters) achieves near-equivalent downstream performance (63.4\% vs.\ 64.0\% for the MLP). This strongly supports our Phase~0 finding: perturbation type is linearly separable in PaliGemma layer~6 representations. The bottleneck is not routing quality---it is expert specialization.

%---------------------------------------------------------------------
\subsection{6-Class vs 3-Class Routing}
\label{sec:6class}

Doubling the number of experts from 3 to 6 (separating visual subtypes) increases parameters from 12.1M to 23.2M but yields only +0.5$\pp$ (64.0\%$\to$64.5\%). Marginal gain per additional expert is 0.17$\pp$. Finer-grained routing offers diminishing returns when the experts themselves are insufficiently specialized.

\begin{table}[h]
\centering
\caption{6-class vs 3-class routing. Nearly doubling parameters yields negligible gain.}
\label{tab:6class}
\small
\begin{tabular}{lcccc}
\toprule
\textbf{Config} & \textbf{Experts} & \textbf{Total Params} & \textbf{Overall (\%)} & $\boldsymbol{\Delta}$ \textbf{vs 3-class} \\
\midrule
3-class routing & 3 & 12.1M & 64.0 $\pm$ 1.6 & --- \\
6-class routing & 6 & 23.2M & 64.5 $\pm$ 1.7 & +0.5 \\
Universal r48 & 1 & 12.0M & 63.3 $\pm$ 1.9 & --- \\
\bottomrule
\end{tabular}
\end{table}

%---------------------------------------------------------------------
\subsection{Rank Sweep: Capacity vs Routing Benefit}
\label{sec:rank_sweep}

We sweep LoRA rank across 16/32/64 and measure routing benefit (routed minus universal) at each rank (Table~\ref{tab:rank_sweep}).

\begin{table}[h]
\centering
\caption{Rank sweep. Routing benefit shows a weak upward trend but is not significant at any tested rank.}
\label{tab:rank_sweep}
\small
\begin{tabular}{lccccc}
\toprule
\textbf{Rank} & \textbf{Universal (\%)} & \textbf{Routed (\%)} & \textbf{Benefit ($\pp$)} & \textbf{$p$-value} & \textbf{Routed Params} \\
\midrule
$r=16$ & 58.5 $\pm$ 1.7 & 59.5 $\pm$ 1.8 & +1.0 & 0.38 & 11.6M \\
$r=32$ & 61.0 $\pm$ 1.8 & 62.2 $\pm$ 1.9 & +1.2 & 0.29 & 23.6M \\
$r=64$ & 65.5 $\pm$ 1.6 & 67.5 $\pm$ 1.8 & +2.0 & 0.14 & 47.8M \\
\bottomrule
\end{tabular}
\end{table}

The routing benefit increases from +1.0$\pp$ at rank~16 to +2.0$\pp$ at rank~64, but none reach statistical significance ($p_{\min} = 0.14$). This weak upward trend is consistent with our oracle-gap analysis: routing may become beneficial at higher expert capacity, but we cannot confirm a capacity threshold within our experimental range. Notably, universal r64 (65.5\%) already exceeds \ours{} at rank~16 (64.0\%), further confirming that capacity scaling dominates routing.

%---------------------------------------------------------------------
\subsection{Layer Placement Ablation}
\label{sec:layer_placement}

We ablate which PaliGemma layers receive LoRA adaptation for each expert type (Table~\ref{tab:layer_placement}).

\begin{table}[h]
\centering
\caption{Layer placement ablation (success rate \%). Best layer subset matches the Phase~0 probing depth ordering but underperforms all-layer placement by 3.8$\pp$ on average.}
\label{tab:layer_placement}
\small
\begin{tabular}{lccccc}
\toprule
\textbf{Expert} & \textbf{Early (2--6)} & \textbf{Mid (6--12)} & \textbf{Late (12--18)} & \textbf{All Layers} & \textbf{Best Subset} \\
\midrule
LoRA-V & 65.5 & 63.0 & 60.2 & 67.0 & Early \\
LoRA-S & 57.0 & 60.5 & 58.0 & 62.0 & Mid \\
LoRA-P & 48.5 & 50.0 & 52.0 & 55.5 & Late \\
\bottomrule
\end{tabular}
\end{table}

The layer preferences confirm the Phase~0 depth ordering: LoRA-V performs best with early layers (65.5\%), LoRA-S with middle layers (60.5\%), and LoRA-P with late layers (52.0\%). This provides mechanistic support for the hypothesis that visual perturbations disrupt early feature extraction, spatial perturbations affect mid-level spatial reasoning, and physical perturbations impact late-stage action planning.

However, all layer-subset placements underperform all-layer adaptation by an average of 3.8$\pp$, indicating that the capacity benefit of full-layer coverage outweighs the precision benefit of layer-targeted placement. This mirrors the broader finding of this paper: perturbation structure is real and detectable, but current LoRA experts cannot exploit it effectively.

%---------------------------------------------------------------------
\subsection{Clean LIBERO: Nominal Performance Degradation}
\label{sec:clean}

Table~\ref{tab:clean} evaluates all methods on unperturbed LIBERO tasks.

\begin{table}[h]
\centering
\caption{Clean LIBERO evaluation. \ours{} introduces a significant $-4.2\pp$ degradation ($p=0.002$) on unperturbed tasks, whereas universal r48 degrades by only $-1.4\pp$ (n.s.).}
\label{tab:clean}
\small
\begin{tabular}{lccc}
\toprule
\textbf{Method} & \textbf{Clean LIBERO (\%)} & $\boldsymbol{\Delta}$ \textbf{vs Frozen ($\pp$)} & \textbf{$p$-value} \\
\midrule
Frozen pi0 & 91.2 $\pm$ 0.8 & --- & --- \\
Universal r16 & 90.5 $\pm$ 0.9 & $-$0.7 & 0.31 \\
Universal r48 & 89.8 $\pm$ 1.0 & $-$1.4 & 0.08 \\
Random routing & 88.2 $\pm$ 1.2 & $-$3.0 & 0.01 \\
\ours{} & 87.0 $\pm$ 1.1 & $-$4.2 & 0.002 \\
\bottomrule
\end{tabular}
\end{table}

\ours{} degrades clean-task performance by 4.2$\pp$ ($p = 0.002$), significantly more than universal r48 ($-1.4\pp$, $p = 0.08$). The router, trained without a ``clean'' category, classifies unperturbed observations as one of the three perturbation types and applies unnecessary LoRA corrections. Random routing ($-3.0\pp$) causes less degradation than learned routing ($-4.2\pp$), suggesting that PAR's confident (but erroneous) classification into a specific perturbation type is more harmful than random dispersal. This is a practical deployment concern: any routing-based system must include an open-set rejection mechanism (e.g., entropy-based gating) to avoid degrading nominal performance.

%---------------------------------------------------------------------
\subsection{Cross-Backbone Validation}
\label{sec:cross_backbone}

We replicate the core pipeline on OpenVLA to test generalizability (Table~\ref{tab:cross_backbone}).

\begin{table}[h]
\centering
\caption{Cross-backbone validation. Phase~0 diagnostic finding replicates on OpenVLA; routing remains non-beneficial.}
\label{tab:cross_backbone}
\small
\begin{tabular}{lccc}
\toprule
\textbf{Metric} & \textbf{pi0 (PaliGemma)} & \textbf{OpenVLA} & \textbf{Consistent?} \\
\midrule
Phase 0: 3-class F1 & 0.993 & 0.952 & Yes (both $>$0.95) \\
Best probing layer & Layer 6/18 & Layer 8/24 & Yes ($\sim$1/3 depth) \\
K=1 F1 & 0.898 & 0.910 & Yes \\
Routing benefit ($\pp$) & +0.7 (vs r48) & +1.6 (vs r48) & Yes (both $<$2$\pp$, n.s.) \\
V$\to$early, S$\to$mid, P$\to$late & Supported & Partial & Partial \\
\bottomrule
\end{tabular}
\end{table}

The Phase~0 finding generalizes: OpenVLA encodes perturbation type with F1 = 0.952, with the optimal probing layer at approximately one-third network depth (layer 8/24 vs.\ layer 6/18 for pi0). The depth ordering is partially replicated (V$\to$early and P$\to$late confirmed; S$\to$mid less clear). The routing benefit on OpenVLA (+1.6$\pp$ vs.\ rank-48 universal) mirrors the pi0 result---positive but not significant. This cross-backbone consistency strengthens the diagnostic finding while confirming that routing failure is not architecture-specific.

%=====================================================================
\section{Discussion}
\label{sec:discussion}

\paragraph{Why does diagnosis succeed but routing fail?}
The PAR router achieves 97.8\% classification accuracy, yet routing provides no significant benefit over universal adaptation. The explanation lies in the expert side: cross-type drops of only 4--6$\pp$ (Table~\ref{tab:cross_drop}) show that rank-16 LoRA experts learn overlapping corrections. The router correctly identifies perturbation type, but the experts do not provide sufficiently differentiated fixes. The soft routing weights confirm this: end-to-end optimization converges to a soft ensemble (Table~\ref{tab:routing_weights}) because mixing experts outperforms strict selection when corrections are interchangeable.

\paragraph{The oracle gap as an open problem.}
Oracle routing's 5.9$\pp$ advantage over learned routing demonstrates that type-specific correction \emph{can} help---under the idealized condition of hard selection with perfectly specialized experts. The rank-sweep trend (Table~\ref{tab:rank_sweep}) suggests routing benefit may grow with expert capacity, though we cannot confirm a crossover point within our experimental range. Future work should investigate whether alternative expert training strategies (e.g., contrastive training, explicit negative perturbation-type regularization) can produce the specialization that current LoRA training does not achieve.

\paragraph{Practical implications.}
For practitioners seeking VLA robustness, our results suggest a simpler recipe: apply a single universal LoRA with the largest affordable rank. Routing adds engineering complexity (router training, soft-merge inference, open-set handling for clean inputs) without measurable benefit under current methods. However, the diagnostic probe itself is cheap ($\sim$0.02M for a linear classifier) and provides valuable runtime monitoring---even without routing, knowing \emph{which} perturbation type is affecting the system has operational value for logging, alerting, and human-in-the-loop intervention.

%=====================================================================
\section{Conclusion}
\label{sec:conclusion}

We presented two findings and one negative result concerning VLA robustness under environmental perturbations. First, VLA internal representations encode perturbation type with near-perfect accuracy at a single layer and timestep, following a consistent depth ordering across architectures. Second, LoRA adaptation of any form substantially improves robustness over frozen baselines (+9.7 to +10.4$\pp$). Third, perturbation-type-specific routing does not outperform parameter-matched universal adaptation: \ours{} (64.0\%) is statistically indistinguishable from a rank-48 universal LoRA (63.3\%, $p=0.42$), and random routing captures most of the multi-expert benefit (62.4\%, $p=0.18$). The 5.9$\pp$ oracle gap and weak rank-sweep trend suggest that routing could become valuable with more specialized experts---we frame this as an open problem rather than a solved contribution. Meanwhile, the diagnostic probe itself is a lightweight, architecture-general tool that we recommend as standard practice for VLA robustness monitoring.

%=====================================================================
\bibliographystyle{plainnat}
\begin{thebibliography}{10}

\bibitem{pi0}
Black, K. et al.
\newblock $\pi_0$: A Vision-Language-Action Flow Model for General Robot Control.
\newblock {\em arXiv preprint arXiv:2410.24164}, 2024.

\bibitem{liberoplus}
Memmel, M. et al.
\newblock LIBERO-Plus: Benchmarking VLA Robustness under Environmental Perturbations.
\newblock In {\em CVPR}, 2026.

\bibitem{vlagse}
Liu, J. et al.
\newblock VLA-GSE: Generalizable and Scalable Expert Routing for VLAs.
\newblock {\em arXiv preprint}, 2025.

\bibitem{ftm}
Memmel, M. et al.
\newblock FTM: Visual Robustness via Fine-Tuning with Masking.
\newblock {\em arXiv preprint}, 2025.

\bibitem{safe}
Zhang, Y. et al.
\newblock SAFE: Perturbation Detection in VLA Policies.
\newblock In {\em RSS}, 2025.

\bibitem{robustvla}
Chen, L. et al.
\newblock RobustVLA: Adversarial Training for VLA Robustness.
\newblock {\em arXiv preprint}, 2025.

\bibitem{lora}
Hu, E. et al.
\newblock LoRA: Low-Rank Adaptation of Large Language Models.
\newblock In {\em ICLR}, 2022.

\bibitem{moira}
Yang, L. et al.
\newblock MoIRA: Mixture of LoRA Experts for Multi-Task Robot Learning.
\newblock {\em arXiv preprint}, 2025.

\bibitem{openvla}
Kim, M. et al.
\newblock OpenVLA: An Open-Source Vision-Language-Action Model.
\newblock In {\em CoRL}, 2024.

\bibitem{peft}
Mangrulkar, S. et al.
\newblock PEFT: Parameter-Efficient Fine-Tuning Methods.
\newblock {\em GitHub repository}, 2023.

\end{thebibliography}

\end{document}
